% Options for packages loaded elsewhere
\PassOptionsToPackage{unicode}{hyperref}
\PassOptionsToPackage{hyphens}{url}
\documentclass[
]{article}
\usepackage{xcolor}
\usepackage{amsmath,amssymb}
\setcounter{secnumdepth}{-\maxdimen} % remove section numbering
\usepackage{iftex}
\ifPDFTeX
  \usepackage[T1]{fontenc}
  \usepackage[utf8]{inputenc}
  \usepackage{textcomp} % provide euro and other symbols
\else % if luatex or xetex
  \usepackage{unicode-math} % this also loads fontspec
  \defaultfontfeatures{Scale=MatchLowercase}
  \defaultfontfeatures[\rmfamily]{Ligatures=TeX,Scale=1}
\fi
\usepackage{lmodern}
\ifPDFTeX\else
  % xetex/luatex font selection
\fi
% Use upquote if available, for straight quotes in verbatim environments
\IfFileExists{upquote.sty}{\usepackage{upquote}}{}
\IfFileExists{microtype.sty}{% use microtype if available
  \usepackage[]{microtype}
  \UseMicrotypeSet[protrusion]{basicmath} % disable protrusion for tt fonts
}{}
\makeatletter
\@ifundefined{KOMAClassName}{% if non-KOMA class
  \IfFileExists{parskip.sty}{%
    \usepackage{parskip}
  }{% else
    \setlength{\parindent}{0pt}
    \setlength{\parskip}{6pt plus 2pt minus 1pt}}
}{% if KOMA class
  \KOMAoptions{parskip=half}}
\makeatother
\setlength{\emergencystretch}{3em} % prevent overfull lines
\providecommand{\tightlist}{%
  \setlength{\itemsep}{0pt}\setlength{\parskip}{0pt}}
\usepackage{bookmark}
\IfFileExists{xurl.sty}{\usepackage{xurl}}{} % add URL line breaks if available
\urlstyle{same}
% fallback for those not using the hyperref driver hyperxmp:
\makeatletter
\@ifundefined{xmpquote}{\newcommand{\xmpquote}[1]{#1}}{}
\makeatother
\hypersetup{
  pdftitle={Java Testing},
  hidelinks,
  pdfcreator={LaTeX via pandoc}}

\title{\texttt{Java} Testing}
\author{}
\date{}

\begin{document}
\maketitle

\subsection{Description}\label{description}

This course covers the full spectrum of testing in the \texttt{Java}
ecosystem. From \texttt{JUnit\ 5} architecture and assertions through
mocking with \texttt{Mockito}, integration testing with
\texttt{Spring\ Boot} Test and \texttt{TestContainers}, to code coverage
analysis with JaCoCo and mutation testing. Students will learn how to
write effective, maintainable tests and adopt \texttt{TDD} as a
development practice.

\subsection{Duration}\label{duration}

16 hours / 2 days

\subsection{Intended Audience}\label{intended-audience}

\begin{itemize}
\tightlist
\item
  \texttt{Java} developers who want to improve their testing skills and
  code quality
\item
  Software engineers adopting test-driven development practices
\item
  Team leads seeking to establish testing standards in their teams
\end{itemize}

\subsection{Prerequisites}\label{prerequisites}

\begin{itemize}
\tightlist
\item
  Intermediate \texttt{Java} programming skills (classes, interfaces,
  generics, annotations)
\item
  Basic understanding of object-oriented design principles
\item
  Familiarity with a build tool such as Maven or \texttt{Gradle}
\item
  Experience with an IDE such as \texttt{IntelliJ\ IDEA} or Eclipse
\end{itemize}

\subsection{Objectives}\label{objectives}

\begin{itemize}
\tightlist
\item
  Master the \texttt{JUnit\ 5} architecture and its extension model
\item
  Write clear, maintainable unit tests with expressive assertions
\item
  Use \texttt{Mockito} to isolate units under test with test doubles
\item
  Implement integration tests using \texttt{Spring\ Boot} Test and
  \texttt{TestContainers}
\item
  Measure and analyze test coverage with JaCoCo
\item
  Apply \texttt{TDD} workflow to drive design and implementation
\item
  Understand and apply mutation testing to evaluate test suite quality
\end{itemize}

\subsection{Outline}\label{outline}

\begin{itemize}
\tightlist
\item
  Introduction to \texttt{Java} Testing

  \begin{itemize}
  \tightlist
  \item
    Why testing matters
  \item
    Types of tests: unit, integration, end-to-end
  \item
    The test pyramid
  \item
    Testing terminology and conventions
  \end{itemize}
\item
  \texttt{JUnit\ 5} Architecture

  \begin{itemize}
  \tightlist
  \item
    \texttt{JUnit} Platform, Jupiter, and Vintage
  \item
    The Jupiter extension model
  \item
    Test classes and methods
  \item
    Display names and nested tests
  \item
    Conditional test execution
  \item
    Test execution order
  \end{itemize}
\item
  Assertions

  \begin{itemize}
  \tightlist
  \item
    Basic assertions: assertEquals, assertTrue, assertNull
  \item
    Grouped assertions and assertAll
  \item
    Exception assertions with assertThrows
  \item
    Timeout assertions
  \item
    Third-party assertion libraries: \texttt{AssertJ} and
    \texttt{Hamcrest}
  \item
    Writing custom matchers
  \end{itemize}
\item
  Parameterized Tests

  \begin{itemize}
  \tightlist
  \item
    @ParameterizedTest annotation
  \item
    Value sources: @ValueSource, @EnumSource, @NullAndEmptySource
  \item
    @CsvSource and @CsvFileSource
  \item
    @MethodSource and @ArgumentsSource
  \item
    Custom argument converters and aggregators
  \end{itemize}
\item
  Test Lifecycle

  \begin{itemize}
  \tightlist
  \item
    @BeforeEach, @AfterEach, @BeforeAll, @AfterAll
  \item
    Test instance lifecycle modes
  \item
    Extension lifecycle callbacks
  \item
    Temporary directory extension
  \item
    Resource management in tests
  \end{itemize}
\item
  \texttt{Mockito} and Test Doubles

  \begin{itemize}
  \tightlist
  \item
    Types of test doubles: stubs, mocks, spies, fakes
  \item
    Creating mocks with \texttt{Mockito}
  \item
    Stubbing method calls
  \item
    Argument matchers and captors
  \item
    Verifying interactions
  \item
    @Mock, @InjectMocks, and @Spy annotations
  \item
    Mocking \texttt{static-methods} and constructors
  \item
    \texttt{BDD\ Mockito} style
  \end{itemize}
\item
  Integration Testing

  \begin{itemize}
  \tightlist
  \item
    Integration test strategies
  \item
    Database testing with embedded databases
  \item
    \texttt{REST\ API} testing with MockMvc
  \item
    WebTestClient for reactive endpoints
  \item
    Testing with real \texttt{HTTP} clients
  \end{itemize}
\item
  \texttt{Spring\ Boot} Testing

  \begin{itemize}
  \tightlist
  \item
    @SpringBootTest and test slices
  \item
    @DataJpaTest for repository testing
  \item
    @WebMvcTest for controller testing
  \item
    @MockBean and @SpyBean
  \item
    Test configuration and profiles
  \item
    Application context caching
  \end{itemize}
\item
  \texttt{TestContainers}

  \begin{itemize}
  \tightlist
  \item
    Introduction to \texttt{TestContainers}
  \item
    Setting up containerized databases for tests
  \item
    \texttt{TestContainers} with \texttt{JUnit\ 5}
  \item
    Custom container definitions
  \item
    \texttt{TestContainers} with \texttt{Spring\ Boot}
  \item
    Testing with \texttt{Kafka}, \texttt{Redis}, and other services
  \end{itemize}
\item
  Test Coverage with JaCoCo

  \begin{itemize}
  \tightlist
  \item
    Understanding code coverage metrics
  \item
    Line, branch, and instruction coverage
  \item
    Configuring JaCoCo with Maven and \texttt{Gradle}
  \item
    Generating and interpreting coverage reports
  \item
    Setting coverage thresholds
  \item
    Excluding classes and packages from coverage
  \end{itemize}
\item
  \texttt{TDD} Workflow

  \begin{itemize}
  \tightlist
  \item
    Red-Green-Refactor cycle
  \item
    Writing the test first
  \item
    Incremental design with \texttt{TDD}
  \item
    \texttt{TDD} and \texttt{clean\ architecture}
  \item
    \texttt{When} \texttt{TDD} works well and when to adapt
  \end{itemize}
\item
  Mutation Testing

  \begin{itemize}
  \tightlist
  \item
    What is mutation testing
  \item
    Mutation operators and mutants
  \item
    PIT mutation testing framework
  \item
    Configuring and running PIT
  \item
    Interpreting mutation testing reports
  \item
    Improving test suites based on mutation results
  \end{itemize}
\end{itemize}

\subsection{Copyright}\label{copyright}

Mark Veltzer
\href{mailto:mark.veltzer@gmail.com}{\nolinkurl{mark.veltzer@gmail.com}},
© 2026

\end{document}
